% ===== §7 Silence as Fulfilment =====
\section{Silence as Fulfilment: The Completion of Three Silences}
\label{sec:silence}

\noindent
\emph{The chapter's distinctive load (demarcated first, to forestall repetition).} \xref{sub:gen-saying} already noted ``silence $=$ no debt, no remainder''; \xref{sec:modes} touched it. This chapter does \emph{not} repeat that. Its distinctive value lies in two things the foregoing did not do: (1)~to gather the three forms ``silence'' takes across this series into a single threefold, making this chapter the \emph{point of convergence} of the silence-line running from Paper~XIII through Paper~XVIII to the present (not one more local claim within Paper~XIX); and (2)~to re-place silence, by way of \xref{sub:gen-hermeneutic}, as a \emph{determinate state of the symbolic}---not the absence of signifiers, but the movement of signifiers having reached a fullness that need say no more. The latter the foregoing did not touch at all.

\subsection{Silence is not aphasia: a re-placement at the symbolic level}
\label{sub:silence-notaphasia}

A fundamental misreading must be barred: the silence of fulfilment is \emph{not} aphasia, \emph{not} a pre-symbolic chaos, \emph{not} the aphasia of the annexed party of \xref{sub:coex-annexation}. Following \xref{sub:gen-hermeneutic} (happiness is conferral of meaning in the symbolic): if silence is a fulfilment, it must be a state \emph{within} the symbolic---the order of signifiers present, the conferral of meaning accomplished, only arrived at a \emph{point of saturation needing no further signifier}.

The decisive contrast:
\begin{itemize}
  \item the aphasia of the annexed (\xref{sub:coex-annexation}) $=$ wanting to speak and having no signifier of one's own to speak with (the symbolic place usurped);
  \item the silence of fulfilment (this chapter) $=$ having signifiers and needing none---the symbolic movement complete, no gap-driven pressure compelling speech.
\end{itemize}
Both ``do not speak''; their standing in the symbolic is opposite---one more pair of ``same appearance, opposite being'' (isomorphic with the two zeros of \xref{sec:dynamics}, the two plainnesses of \xref{sub:neuro-fullness}).

\subsection{The threefold of three silences (load-bearing; the series' convergence)}
\label{sub:silence-three}

\begin{itemize}
  \item \emph{Silence one---obligation (Paper XIII \citep{paperxiii}).} Generative silence $=$ \emph{not to occupy} the other's generative space. An \emph{ought}---to yield place actively, that the other may generate of itself. Directed at: the other's freedom to generate.
  \item \emph{Silence two---restraint (Paper XVIII \citep{paperxvii}).} The negative structural obligation $=$ \emph{not to dissolve} the other's generativity (not to crush the other with one's own order). Also an \emph{ought}---to restrain the tendency to annex. Directed at: the preservation of the other's subjectivity (cf.\ the place of conferral, \xref{sub:coex-annexation}).
  \item \emph{Silence three---needlessness (Paper XIX, this paper).} The silence of fulfilment $=$ \emph{nothing to say and no need to say it}. Not an ought-not-to-speak, but the absence of debt, of remainder, of any unresolved tension compelling speech. Directed at: the fulfilment of the relation itself (the gap abiding in the mode of fullness, \xref{sub:gen-twomodes}).
\end{itemize}

\emph{The structure of the ascent.} The first two are \emph{negative oughts} (``one ought not to \ldots'')---they remain obligations, still presupposing a tendency to be checked (occupation, annexation). The third is a \emph{fulfilled actuality} (``there is no need to \ldots'')---no longer an obligation, because there is no longer a tendency to check; it is the state \emph{after} obligation has been wholly satisfied, in which obligation itself dissolves. The three are therefore not coordinate but an \emph{ascent}: from ``ought not to occupy'' (XIII) to ``ought not to dissolve'' (XVIII) to ``there is no need to speak'' (XIX)---obligation \emph{sublated} (\emph{Aufhebung}: at once preserved and surpassed) in fulfilment.

\begin{claim}[The three silences]
Silence completes a series. The silence of obligation does not occupy; the silence of restraint does not dissolve; the silence of needlessness has nothing to say and no need to say it. The first two are negative oughts that presuppose a tendency to be checked; the third is the fulfilled state in which, no tendency remaining, the obligation itself dissolves. Silence is thus not the absence of meaning but the saturation of it: signifiers present, and none required.
\end{claim}

\subsection{The Eastern interface: arriving from within, not borrowing authority}
\label{sub:silence-east}

Daoism: \emph{tiandan} (the plain), the ``savouring of the flavourless'' (\emph{wei wu wei}), ``the great sound is faint'' (\emph{dayin xisheng}), ``one who knows does not speak'' \citep{laozi1963,zhuangzi2009}. ``Savouring the flavourless'' fits the third silence especially: not the absence of flavour (aphasia, flavourlessness) but flavour grown so complete it need add no more. The Chan ``no establishment of words'' (\emph{buli wenzi}) is \emph{not} ``no words'' but the needlessness of further words.

The voice-discipline (per the paper throughout): these are not authorities adduced to stamp the argument; the paper has \emph{arrived again, from within} the argument of the generative relation, at the place they once reached. The convergence is corroboration, not borrowed light.

\subsection{Close and transition}
\label{sub:silence-close}

The threefold is complete: obligation (yielding) $\to$ restraint (not dissolving) $\to$ needlessness (fulfilment). The third silence is the signature, at the level of saying, of this paper's happiness---fulfilment need not be stated; it need only be lived, a circle at a time, and kept, together, in silence.

Transition to \xref{sec:unwritability}: yet this paper \emph{is} a paper; it is, right now, saying the very thing that ``needs no saying.'' This performative tension is no error---the next chapter shows it to be the self-verification of the thesis.
